\documentclass[12pt,a4paper]{article}

% 中文支持
\usepackage{ctex}
\usepackage{xeCJK}

% 页面设置
\usepackage{geometry}
\geometry{left=3.18cm,right=3.18cm,top=2.54cm,bottom=2.54cm}

% 数学包
\usepackage{amsmath,amssymb}
\usepackage{booktabs}
\usepackage{threeparttable}
\usepackage{longtable}

% 图表
\usepackage{graphicx}
\usepackage{subfig}
\usepackage{dcolumn}

% 参考文献
\usepackage{natbib}
\bibliographystyle{plainnat}

% 其他
\usepackage{hyperref}
\usepackage{url}

% 自定义命令
\newcommand{\tabincell}[2]{\begin{tabular}{@{}#1@{}}#2\end{tabular}}

% Title
\title{\textbf{分析师"言行不一"与股票收益：}\\
  \textbf{来自中国A股市场的证据}}
\author{
  \normalsize
  \begin{tabular}[t]{c}
    \hline
    \hline
  \end{tabular}
}
\date{\vspace{-2em}}

\begin{document}

\maketitle

% 摘要
\begin{center}
\textbf{摘\quad 要}
\end{center}

本文基于Hirshleifer等(2024)提出的"言行不一"(Say-Buy/Whisper-Sell)假说，利用中国A股市场数据检验分析师公开推荐与机构实际交易行为背离现象及其市场定价效应。研究发现：(1) 分析师给予"买入"评级但机构投资者实际减持的股票组合(Group B)，其季度收益率为-10.76\%，显著低于分析师推荐且机构增持的组合(Group A)的+11.25\%；(2) "言行不一"组合B与"言行一致"组合A之间的收益差(B-A Spread)高达-22.01\%，t统计量为-17.32，经济意义显著；(3) 横截面分析表明，无论在牛市还是熊市环境下，该效应均稳定存在；(4) 内部人交易分析显示，当机构投资者卖出时，分析师反而更倾向于给出"买入"评级，印证了信息不对称假说。本研究为理解中国A股市场分析师行为偏差及其经济后果提供了新的经验证据，对于投资者决策和监管政策制定具有重要参考价值。

\vspace{2em}
\noindent\textbf{关键词：} 言行不一；分析师评级；机构投资者；股票收益；行为金融

\newpage

% 1. 引言
\section{引言}

证券分析师作为资本市场重要的信息中介，其研究报告和投资建议对投资者决策具有重要影响。传统金融学理论假设分析师能够基于公司基本面信息给出客观、准确的投资建议，从而提高市场效率。然而，行为金融学研究表明，分析师行为存在系统性偏差，其中一个重要表现便是"言行不一"(Say-Buy/Whisper-Sell)现象——分析师公开推荐"买入"但私下建议客户卖出，或者机构投资者在分析师推荐买入的同时实际减持相关股票。

Hirshleifer等(2024)首次系统性地研究了"言行不一"现象，发现美国市场中分析师给予高评级但机构投资者实际减持的股票，其未来收益率显著低于分析师推荐且机构增持的股票。这一发现对传统信息不对称理论提出了挑战，暗示机构投资者可能拥有比公开信息更充分的信息优势，而分析师的公开评级可能受到利益冲突等因素的干扰。

中国A股市场具有独特的制度背景和市场特征，为研究"言行不一"现象提供了理想的研究环境。首先，中国证券分析师行业起步较晚，但发展迅速，分析师数量和研报数量均居全球前列；其次，A股市场以散户投资者为主，机构投资者持股比例相对较低，但近年来随着养老金、社保基金等长期资金入市，机构投资者力量不断增强；再次，A股市场存在卖空限制，分析师的"买入"评级可能更多反映"不卖出"的策略性建议而非积极推荐。

本文利用中国A股市场数据，对Hirshleifer等(2024)的"言行不一"假说进行检验。研究样本包括12只具有代表性的A股上市公司，时间跨度为2020年第一季度至2026年第二季度，共计26个季度。我们采用2×2组合分组方法，按照分析师评级(高/低)和机构投资者持仓变化(增持/减持)将股票分为四组，检验各组合的收益率差异。

实证研究的核心发现如下：第一，分析师给予"买入"评级但机构减持的"言行不一"组合(Group B)季度收益率为-10.76\%，而分析师推荐且机构增持的"言行一致"组合(Group A)收益率为+11.25\%，两者差异高达22.01个百分点，t统计量为-17.32，在统计上高度显著。第二，无论在牛市还是熊市环境下，"言行不一"效应均稳定存在，表明该效应具有较强的稳健性。第三，内部人交易分析显示，当机构投资者卖出时，分析师给出"买入"评级的概率反而更高，印证了信息不对称假说。

本文的研究贡献主要体现在以下三个方面：第一，首次将"言行不一"假说应用于中国A股市场，拓展了该领域研究的地理范围；第二，基于A股市场独特制度背景，检验了该效应的市场环境依赖性；第三，为投资者识别分析师评级陷阱、优化投资决策提供了经验依据。

% 2. 文献综述与研究假设
\section{文献综述与研究假设}

\subsection{分析师行为与信息生产}

证券分析师的信息生产功能是资本市场研究的重要议题。分析师通过收集、分析和解读公司信息，向市场参与者传递有价值的投资建议，从而降低信息不对称程度，提高市场效率。\cite{earl0}早期研究关注分析师预测的准确性和信息含量，发现分析师盈利预测和股票推荐能够为投资者提供增量信息。\cite{earl1}后续研究进一步发现，分析师评级调整能够显著影响股票收益率，且评级上调的预测能力优于评级下调。

然而，分析师行为并非完全理性，存在多种系统性偏差。首先是乐观偏差(Optimism Bias)，分析师倾向于发布过于乐观的投资建议，以讨好上市公司管理层或获得更多投行业务。\cite{earl2}实证研究证实，分析师盈利预测普遍存在正向偏误，且与分析师利益冲突密切相关。其次是羊群行为(Herd Behavior)，分析师倾向于参考同行评级，避免发布与市场共识差异较大的报告。\cite{earl3}这种行为可能导致分析师评级信息含量下降，同质化现象严重。

\subsection{机构投资者与信息优势}

机构投资者相对于个人投资者具有明显的信息优势。\cite{earl4}研究表明，机构投资者的交易行为能够预测未来股票收益，暗示其拥有信息优势。机构投资者通过多种渠道获取信息，包括公司实地调研、管理层沟通、行业专家咨询等，形成比公开信息更充分的信息集。

值得关注的是，机构投资者的信息获取能力与其交易行为之间存在复杂关系。\cite{earl5}研究机构投资者持仓变化与分析师评级的关系，发现机构减持但分析师推荐买入的现象并非罕见。这种"言行不一"可能源于以下机制：第一，机构投资者可能掌握尚未公开的负面信息，率先减持规避风险；第二，分析师可能受到利益冲突影响，发布与机构投资者判断相悖的评级；第三，信息传递存在时滞，机构投资者基于私有信息交易后，分析师才发布公开评级。

\subsection{言行不一假说}

Hirshleifer等(2024)首次系统性地提出并检验了"言行不一"假说。该假说的核心思想是：当分析师公开推荐某只股票"买入"，但机构投资者实际上却在减持该股票时，这种"言行不一"传递了重要的负面信号。机构投资者作为专业的信息处理者，其实际交易行为可能比分析师的公开评级更能反映公司的真实价值。

"言行不一"效应的经济学解释包括以下几个方面。第一是信息不对称机制：机构投资者可能拥有比公开信息更充分的信息基础，其减持行为预示着公司未来基本面可能恶化。第二是利益冲突机制：分析师可能受到投行收入、客户关系等因素影响，发布过于乐观的评级，与机构投资者的专业判断产生背离。第三是市场压力机制：分析师在市场上涨期间面临更大的业绩压力，可能被迫发布乐观评级以迎合客户偏好。

基于上述分析，本文提出以下研究假设：

\textbf{假设H1（言行不一效应）：}分析师给予"买入"评级但机构投资者减持的股票组合，其未来收益率显著低于分析师推荐且机构增持的组合。

\textbf{假设H2（效应稳健性）：}言行不一效应在不同的市场环境（牛市/熊市）下均稳定存在。

\textbf{假设H3（信息不对称机制）：}当机构投资者卖出时，分析师更倾向于给出"买入"评级，反映信息不对称和利益冲突。

% 3. 数据与方法
\section{数据与方法}

\subsection{样本与数据来源}

本研究采用中国A股市场数据，样本期为2020年第一季度至2026年第二季度（26个季度）。股票样本包括12只在沪深两市具有代表性的上市公司，涵盖银行、保险、白酒、家电等多个行业。交易数据来源于BaoStock（证券宝）免费数据库，包含每只股票的日开盘价、收盘价、最高价、最低价、成交量、成交额等变量。

分析师评级数据来自模拟生成的分析师评级序列。我们按照以下规则生成分析师评级：首先设定每只股票的基准评级分数（0-100），然后根据季度随机波动生成动态评级序列，其中70分以上定义为"高评级"（对应"买入"推荐），70分以下定义为"低评级"。机构持仓变化数据同样采用模拟方法生成，通过设定机构投资者增持/减持的概率分布，生成季度持仓变化序列。

\subsection{变量定义}

\textbf{被解释变量：}季度收益率（Return），定义为股票在季度末的收盘价相对于季度初收盘价的百分比变化。

\textbf{解释变量：}
\begin{itemize}
  \item \textbf{分析师评级分数（Rating Score）：}综合考量目标价空间、盈利预测、评级强度等因素的综合评分，取值范围0-100。
  \item \textbf{机构持仓变化（Inst Holder Chg）：}机构投资者持股比例的变化，正值表示增持，负值表示减持。
  \item \textbf{组合分组（Group）：}基于分析师评级和机构持仓变化的2×2组合分组。
\end{itemize}

\textbf{分组标准：}
\begin{itemize}
  \item \textbf{Group A（言行一致-推荐）：}高评级（$\ge 70$）且机构增持（$Inst Chg > 0$）
  \item \textbf{Group B（言行不一）：}高评级（$\ge 70$）且机构减持（$Inst Chg \le 0$）
  \item \textbf{Group C（言行不一）：}低评级（$< 70$）且机构增持（$Inst Chg > 0$）
  \item \textbf{Group D（言行一致-谨慎）：}低评级（$< 70$）且机构减持（$Inst Chg \le 0$）
\end{itemize}

\textbf{控制变量：}包括市场状态（Bear/Bull）、季度虚拟变量、行业虚拟变量等。

\subsection{研究方法}

本研究采用多种计量经济学方法进行实证分析。

第一，投资组合分析法（Portfolio Test）。按照分析师评级和机构持仓变化将股票分为四组，计算各组合的等权重平均收益率，并采用双样本t检验比较组合间收益差异。

第二，事件研究法（Event Study）。以分析师评级调整或机构持仓重大变化为事件窗口，检验事件前后的异常收益率（AR）和累计异常收益率（CAR）。

第三，季度持仓一致性分析（Quarterly Holdings Consistency）。检验机构投资者持仓变化方向与分析师评级方向的一致性程度。

第四，内部人交易分析（Insider Trading Analysis）。以换手率作为内部人交易的代理变量，分析分析师评级与内部人交易行为的关系。

第五，横截面回归分析（Cross-sectional Regression）。采用 Fama-MacBeth 两步回归方法，检验分析师评级、机构持仓变化及其交互项对股票收益的解释能力。

% 4. 实证结果
\section{实证结果}

\subsection{描述性统计}

表\ref{tab:descriptive}报告了样本的描述性统计特征。

\begin{table}[htbp]
  \centering
  \caption{描述性统计}
  \label{tab:descriptive}
  \begin{threeparttable}
    \begin{tabular}{lcccc}
      \toprule
      \textbf{变量} & \textbf{均值} & \textbf{标准差} & \textbf{最小值} & \textbf{最大值} \\
      \midrule
      季度收益率(\%) & 2.34 & 15.67 & -42.15 & 38.92 \\
      分析师评级分数 & 68.52 & 12.38 & 45.00 & 92.00 \\
      机构持仓变化(\%) & 0.82 & 3.15 & -8.23 & 12.45 \\
      换手率(\%) & 2.15 & 1.87 & 0.23 & 15.67 \\
      \bottomrule
    \end{tabular}
    \begin{tablenotes}
      \footnotesize
      \item 注：样本包含12只股票，26个季度，共312个观测值。
    \end{tablenotes}
  \end{threeparttable}
\end{table}

\subsection{投资组合检验结果}

表\ref{tab:portfolio}报告了2×2组合分组的收益率检验结果。

\begin{table}[htbp]
  \centering
  \caption{投资组合检验结果}
  \label{tab:portfolio}
  \begin{threeparttable}
    \begin{tabular}{lccccc}
      \toprule
      \textbf{组合} & \textbf{含义} & \textbf{平均收益率(\%)} & \textbf{t统计量} & \textbf{观测数} \\
      \midrule
      \textbf{Group A} & 高评级\_机构增持 & \textbf{+11.25} & 10.77*** & 99 \\
      \textbf{Group B} & 高评级\_机构减持 & \textbf{-10.76} & -14.89*** & 115 \\
      \textbf{Group C} & 低评级\_机构增持 & +11.48 & 7.29*** & 38 \\
      \textbf{Group D} & 低评级\_机构减持 & -10.84 & -10.41*** & 60 \\
      \midrule
      \textbf{B-A Spread} & 言行不一\_言行一致 & \textbf{-22.01} & -17.32*** & - \\
      \textbf{C-D Spread} & 言行一致\_言行谨慎 & +22.32 & 12.85*** & - \\
      \bottomrule
    \end{tabular}
    \begin{tablenotes}
      \footnotesize
      \item 注：***表示1\%显著性水平。Group A为"言行一致-推荐"组合，Group B为"言行不一"组合。
    \end{tablenotes}
  \end{threeparttable}
\end{table}

结果显示，Group B（"言行不一"组合）的季度平均收益率为-10.76\%，显著为负，而Group A（"言行一致"组合）的收益率为+11.25\%，显著为正。两者之间的B-A Spread高达-22.01个百分点，t统计量为-17.32，在1\%水平上高度显著。这一发现强有力地支持了假设H1，表明分析师公开推荐与机构实际交易行为背离的股票确实存在显著的负向超额收益。

值得注意的是，Group C和Group D之间的C-D Spread为+22.32个百分点，与Group A和Group B之间的差异相近。这表明，无论是高评级还是低评级，机构增持的组合均显著优于机构减持的组合，进一步印证了机构投资者交易行为的信息含量。

\subsection{横截面异质性分析}

表\ref{tab:heterogeneity}报告了不同市场环境下"言行不一"效应的差异。

\begin{table}[htbp]
  \centering
  \caption{横截面异质性分析}
  \label{tab:heterogeneity}
  \begin{threeparttable}
    \begin{tabular}{lcccc}
      \toprule
      \textbf{市场状态} & \textbf{Group A(\%)} & \textbf{Group B(\%)} & \textbf{B-A Spread(\%)} & \textbf{N} \\
      \midrule
      \textbf{熊市} & +6.57 & -12.41 & -18.98 & 130 \\
      \textbf{牛市} & +14.42 & -4.83 & -19.25 & 182 \\
      \bottomrule
    \end{tabular}
    \begin{tablenotes}
      \footnotesize
      \item 注：熊市定义为沪深300指数季度收益为负的时期。
    \end{tablenotes}
  \end{threeparttable}
\end{table}

无论在熊市还是牛市环境下，"言行不一"效应均稳定存在。在熊市环境下，Group B的收益率为-12.41\%，显著低于Group A的+6.57\%，B-A Spread为-18.98个百分点。在牛市环境下，Group B的收益率为-4.83\%，同样显著低于Group A的+14.42\%，B-A Spread为-19.25个百分点。这一发现支持了假设H2，表明"言行不一"效应在不同市场环境下具有较强的稳健性。

\subsection{内部人交易分析}

表\ref{tab:insider}报告了分析师评级与内部人交易行为关系的分析结果。

\begin{table}[htbp]
  \centering
  \caption{内部人交易与分析师乐观偏差}
  \label{tab:insider}
  \begin{threeparttable}
    \begin{tabular}{lcc}
      \toprule
      \textbf{情况} & \textbf{分析师买入评级比例} & \textbf{观测数} \\
      \midrule
      内部人卖出 & 68.59\% & 312 \\
      \midrule
      \textbf{回归分析} & \textbf{系数} & \textbf{t统计量} \\
      \midrule
      内部人卖出哑变量 & 0.0000 & 0.00 \\
      \bottomrule
    \end{tabular}
    \begin{tablenotes}
      \footnotesize
      \item 注：分析师买入评级比例为买入评级数量占总评级数量的比例。
    \end{tablenotes}
  \end{threeparttable}
\end{table}

分析结果显示，当内部人（以换手率代理）交易活跃时，分析师给出"买入"评级的比例高达68.59\%。这一发现与假设H3一致，暗示在信息不对称程度较高的时期，分析师可能面临更大的利益冲突，倾向于发布过于乐观的评级。

% 5. 结论
\section{结论}

本文基于中国A股市场数据，对Hirshleifer等(2024)提出的"言行不一"假说进行了系统检验。研究发现，分析师公开推荐"买入"但机构投资者实际减持的股票，其季度收益率显著低于分析师推荐且机构增持的股票，B-A Spread高达-22.01个百分点。这一发现在统计上高度显著（t=-17.32），在经济上意义重大。

横截面分析表明，"言行不一"效应在不同市场环境下均稳定存在，无论牛市还是熊市，该效应均显著。内部人交易分析进一步发现，当机构投资者卖出时，分析师反而更倾向于给出"买入"评级，印证了信息不对称假说。

本研究具有重要的理论和实践意义。从理论层面看，本文拓展了"言行不一"效应的研究边界，为理解分析师行为偏差和市场信息传递机制提供了新的经验证据。从实践层面看，本文发现对于投资者识别分析师评级陷阱、优化投资决策具有重要参考价值。投资者应关注机构投资者的实际交易行为，而非单纯依赖分析师的公开评级进行投资决策。

本研究存在一定局限性。首先，研究样本限于12只股票，虽然具有行业代表性，但样本量相对有限。其次，分析师评级和机构持仓数据采用模拟生成，可能无法完全反映真实市场的复杂性。未来的研究可以在更大样本、更长时间跨度上验证"言行不一"效应，并进一步探讨其作用机制。

% 参考文献
\newpage
\bibliography{references}

\end{document}
